\chapter{How to Use This Template}

This short front-matter chapter doubles as a quick reference. It is an
unnumbered chapter (because it sits in \verb|\frontmatter|), so it gets the
same opener as a numbered chapter, minus the ``Chapter'' line, and appears
in bold in the table of contents.

\subsection*{Structure}
\begin{itemize}[leftmargin=1.5em]
  \item \verb|\chapter{Title}| then, optionally,
        \verb|\chapterquote{Quote}{Author}{Source}|.
  \item \verb|\section|, \verb|\subsection| as usual. Sections are numbered
        $1, 2, \dots$ inside each chapter, subsections $1.1, 1.2, \dots$
  \item Theorem-like environments: \texttt{theorem}, \texttt{proposition},
        \texttt{lemma}, \texttt{corollary}, \texttt{claim},
        \texttt{conjecture}, \texttt{definition}, \texttt{example},
        \texttt{notation}, \texttt{remark} (plus starred, unnumbered
        versions). All share one counter that restarts in every section.
        An optional argument is the subtitle:
        \verb|\begin{theorem}[Bolzano--Weierstrass]|.
  \item \texttt{proof} works as in \texttt{amsthm}.
\end{itemize}

\subsection*{References}
Label anything with \verb|\label{...}| and cite it with \verb|\cref{...}|
(or \verb|\Cref| at the start of a sentence). Inside the same chapter you get
``Theorem~1.6'' and ``\S2''; from another chapter the chapter number is
added automatically: ``Theorem~I.1.6'' and ``\S I.2''.

\subsection*{Exercises}
At the end of a section:
\begin{verbatim}
\begin{exercises}
  \begin{exercise}[optional title] \label{ex:foo}
    Prove that ... \fromtext{thm:bar, sec:baz}
  \end{exercise}
  \begin{exercise}
    Show that ... \fromex{ex:qux}
  \end{exercise}
\end{exercises}
\end{verbatim}
\verb|\fromtext{...}| puts a \MBtextmarker{} before the exercise (it is
referred to from the text); \verb|\fromex{...}| puts a \MBexmarker{} (it is
referred to from other exercises). Either way the places you list are printed,
hyperlinked, in brackets at the end of the exercise. Use
\verb|exercise*| to flag a harder exercise with an asterisk.

\subsection*{Problems}
At the end of a chapter, \verb|\begin{problems} ... \end{problems}| starts a
numbered ``Problems'' section. Inside, use \verb|problem| or \verb|problem*|;
\verb|\hint{...}| prints a bracketed hint on its own line.

\subsection*{Pseudocode}
Procedures are typeset in the style of CLRS, with typed parameters:
\begin{verbatim}
\begin{procedure}[returns=\N, columns={cost, times},
                  caption={...}, label=alg:foo, float]
  {Insertion-Sort}{A: \R^n, n: \N}
  \For{$j = 2$ \To $n$}      \Cols{c_1 & n}
    \State $key = A[j]$      \Cols{c_2 & n-1}
  \EndFor
\end{procedure}
\end{verbatim}
Useful options are \verb|float| (let it float), \verb|figure| (number it as
a Figure instead of an Algorithm), \verb|wrap| (set it beside the text),
\verb|width|, \verb|colwidth(s)| and \verb|numbers=false|. Without
\verb|caption| and \verb|float| it sits inline, unnumbered, as in CLRS.
The \texttt{procedure} environment is documented in the class file.

\subsection*{Footnotes}
\verb|\footnote{...}| works as usual. Footnotes are styled as in Aluffi and
numbered afresh in each chapter.
